%% =========================================================
%% Example: the anomalous Z2 x Z2 symmetry with a mixed
%% three-cocycle, its condensation defect, and the sixteen
%% pair blocks of the stacked defect tensor.
%% Covers Notes/OpenProblemsTN/checks/asym_z2z2_typeIII_data.md.
%% =========================================================

\section{The \texorpdfstring{$\Z_2\times\Z_2$}{Z2 x Z2} symmetry example}%
\label{sec:asymex_z2z2}

On a periodic chain of $L$ sites, each carrying two qubits with label
$a=(a_1,a_2)$, four matrix product operators
$U_e=\id$, $U_x$, $U_y$, $U_{xy}$ satisfy
$U_gU_h=\lambda(g,h)^LU_{gh}$ for $L\geq1$, where
$\lambda(g,h)=(-1)^{g_1h_2}$. Here $x=(1,0)$, $y=(0,1)$, and addition
of qubit labels is modulo two. The tensor of $U_g$ is
\begin{align}
    M_g^{(o,i)}[\alpha,\beta]
     & =\delta_{o,g+i} \delta_{\beta,\pi_g(i)}\,f_g(\alpha,i),
    \label{eq:asymex_z2z2_tensor}
\end{align}
where $\pi_g(i)=i_2$ and $f_g(\alpha,i)=(-1)^{\alpha i_2}$ for
$g\in\{x,xy\}$. For $g\in\{e,y\}$, the bond dimension is one,
$\pi_g(i)=0$, and $f_g(\alpha,i)=1$.
Every entry lies in $\{0,\pm1\}$.
The dependence of $M_x$ on the bond indices and the second input
qubit uses the two matrices of the bond-two symmetry tensor in
Section~\ref{sec:asymex_czx}; $M_{xy}$ is obtained from $M_x$ by
relabeling the outgoing physical index as $o\mapsto o+y$.

The results below concern the word spans of these tensors, their
periodic multiplication law, and the compression of the sixteen pair
blocks in the square of their direct sum. A three-cocycle class requires
a separate comparison of the two orders of fusion of three tensors.
Neither the periodic multiplication law nor the absence of a nonzero
sitewise intertwiner alone computes that class.
% The mixed-class interpretation and its restrictions to x, y, and xy
% are external data in Notes/OpenProblemsTN/checks/asym_z2z2_typeIII_data.md,
% Sections 1--3; no three-cocycle class is computed by the declarations below.

\begin{center}
    \tenkzkernel
    % Formula: M_g^{(o,i)}[\alpha,\beta] for g in {e,x,y,xy}, one periodic
    % ring per symmetry element, drawn once with a generic label.
    % Ink/index correspondence: every horizontal virtual index is
    % contracted, including the return bond; the upper legs carry the free
    % outgoing physical indices o_1,...,o_L and the lower legs carry the
    % free incoming physical indices i_1,...,i_L. The skin=dots atom
    % abbreviates the omitted cells, so the drawn and logged physical
    % arities are necessarily schematic rather than L. Boundary signature
    % (virtual-west,virtual-east,physical-up,physical-down) = (0,0,L,L).
    \begin{tenkz}[rows={op}, cols=4, boundary=periodic]
        \tn[skin=mpo, ports={90:physical:$o_1$, 270:physical:$i_1$}]{M_g} &
        \tn[skin=mpo, ports={90:physical:$o_2$, 270:physical:$i_2$}]{M_g} &
        \tn[skin=dots, ports={180:virtual, 0:virtual}]{} &
        \tn[skin=mpo, ports={90:physical:$o_L$, 270:physical:$i_L$}]{M_g}
    \end{tenkz}
\end{center}

\begin{definition}[The four symmetry tensors]%
    \label{def:asymex_z2z2_tensors}
    \lean{Z2Z2Condensation.eTensor, Z2Z2Condensation.yTensor,
        Z2Z2Condensation.xTensor, Z2Z2Condensation.xyTensor,
        Z2Z2Condensation.eMPS, Z2Z2Condensation.yMPS, Z2Z2Condensation.xMPS,
        Z2Z2Condensation.xyMPS}
    \leanok
    \uses{def:mpo_tensor, def:mpo_to_mps}
    The tensors of $U_e$ and $U_y$ have bond dimension one, with entries
    $\delta_{o,i}$ and $\delta_{o,y+i}$, respectively. Each has four
    nonzero physical-index pairs. The tensors of $U_x$ and $U_{xy}$ have
    bond dimension two and the four nonzero matrices
    specified by~(\ref{eq:asymex_z2z2_tensor}), each equal to
    $\begin{pmatrix}1&0\\1&0\end{pmatrix}$ or
    $\begin{pmatrix}0&1\\0&-1\end{pmatrix}$ according to $i_2$.
    Identifying the two-qubit labels with $\{0,1,2,3\}$ and pairing the
    physical indices gives an MPS tensor on the sixteen-letter
    alphabet indexed by $4o+i$.
\end{definition}

\begin{theorem}[Injectivity after blocking of the four symmetry tensors]%
    \label{thm:asymex_z2z2_normal}
    \lean{Z2Z2Condensation.eMPS_isNormal, Z2Z2Condensation.yMPS_isNormal,
        Z2Z2Condensation.xMPS_isNormal, Z2Z2Condensation.xyMPS_isNormal,
        Z2Z2Condensation.xMPS_isNBlkInjective_two,
        Z2Z2Condensation.xyMPS_isNBlkInjective_two,
        Z2Z2Condensation.xMPS_not_isNBlkInjective_one,
        Z2Z2Condensation.xyMPS_not_isNBlkInjective_one}
    \leanok
    \uses{def:asymex_z2z2_tensors, def:normal}
    The bond-one tensors of $U_e$ and $U_y$ are injective without
    blocking: a nonzero matrix is the scalar $1$ and spans
    $\MN{1}$. The bond-two tensors of $U_x$ and $U_{xy}$ are not
    injective without blocking, since their four nonzero matrices span
    only a two-dimensional space. Their length-two words span
    $\MN{2}$, so they are injective after blocking two sites.
    These assertions concern word spans. The tensors are proportional
    to normal tensors (NT); the spectral-radius normalization required
    for NT is imposed separately.
\end{theorem}

\begin{proof}
    \leanok
    For each bond-one tensor, a nonzero scalar matrix spans $\MN{1}$.
    For each bond-two tensor, explicit integer linear combinations of
    four length-two words give the four matrix units. At length one,
    the two distinct nonzero matrices are linearly independent, and
    a nonzero linear functional vanishing on both shows that they do
    not span $\MN{2}$.
\end{proof}

\section{The condensation defect and its periodic fusion table}%
\label{sec:asymex_z2z2_defect}

The \emph{condensation defect} is the sum of the four periodic operators,
$A=U_e+U_x+U_y+U_{xy}$. It is generated by the direct-sum tensor
$N=M_e\oplus M_x\oplus M_y\oplus M_{xy}$ of
Definition~\ref{def:mpo_direct_sum}, with bond dimension six.
In this example, $N$ denotes the direct-sum tensor and $L$ denotes
the chain length.

\begin{center}
    \tenkzkernel
    % Formula: the periodic ring of the direct-sum tensor N, generating
    % A = U_e + U_x + U_y + U_xy by Theorem~\ref{thm:mpo_direct_sum_additive}
    % applied three times.
    % Ink/index correspondence: as in the ring above, with the node
    % renamed N; boundary signature (0,0,L,L).
    \begin{tenkz}[rows={op}, cols=4, boundary=periodic]
        \tn[skin=mpo, ports={90:physical:$o_1$, 270:physical:$i_1$}]{N} &
        \tn[skin=mpo, ports={90:physical:$o_2$, 270:physical:$i_2$}]{N} &
        \tn[skin=dots, ports={180:virtual, 0:virtual}]{} &
        \tn[skin=mpo, ports={90:physical:$o_L$, 270:physical:$i_L$}]{N}
    \end{tenkz}
\end{center}

\begin{definition}[The condensation defect tensor]%
    \label{def:asymex_z2z2_defect}
    \lean{Z2Z2Condensation.condensationTensor,
        Z2Z2Condensation.mpo_condensationTensor}
    \leanok
    \uses{def:mpo_direct_sum, def:asymex_z2z2_tensors}
    The condensation defect tensor is the iterated direct sum
    $N=((M_e\oplus M_x)\oplus M_y)\oplus M_{xy}$ of bond dimension six.
    Its periodic operator is the sum of the four periodic operators of
    Definition~\ref{def:asymex_z2z2_tensors}: $\rho^{(L)}(N)=\rho^{(L)}(M_e)
        +\rho^{(L)}(M_x)+\rho^{(L)}(M_y)+\rho^{(L)}(M_{xy})$ at every system
    size.
\end{definition}

\begin{proof}
    \leanok
    \uses{thm:mpo_direct_sum_additive}
    Apply the additivity of
    Theorem~\ref{thm:mpo_direct_sum_additive} three times.
\end{proof}

\begin{theorem}[The periodic fusion table]%
    \label{thm:asymex_z2z2_fusion}
    \lean{Z2Z2Condensation.bondDim, Z2Z2Condensation.symTensor,
        Z2Z2Condensation.gmul, Z2Z2Condensation.fusionSign,
        Z2Z2Condensation.mpo_symTensor_mul,
        Z2Z2Condensation.mpo_x_mul_y_eq_neg_of_odd}
    \leanok
    \uses{def:asymex_z2z2_tensors, lem:asym_mpo_trace_eq}
    For every $L\geq1$ and every $g,h\in\Z_2\times\Z_2$, the periodic
    operators satisfy
    $\rho^{(L)}(M_g)\rho^{(L)}(M_h)=\lambda(g,h)^L
        \rho^{(L)}(M_{gh})$, where
    \begin{align}
        \lambda(g,h)
         & =\begin{cases}
                -1 & \text{if }g\in\{x,xy\}\text{ and }h\in\{y,xy\}, \\
                +1 & \text{otherwise}.
            \end{cases}
        \notag
    \end{align}
    On odd rings, $\lambda(x,y)\lambda(y,x)=-1$, so
    $\rho^{(L)}(M_x)\rho^{(L)}(M_y)=-\rho^{(L)}(M_y)\rho^{(L)}(M_x)$.
    The four operators form a projective representation of
    $\Z_2\times\Z_2$ on odd rings and a linear representation on even
    rings.
\end{theorem}

\begin{proof}
    \leanok
    \uses{lem:asym_mpo_trace_eq}
    For each of the sixteen pairs, a finite calculation with the integer
    matrices of the stacked tensor gives
    $\tr((M_g\cdot M_h)^w)=\lambda(g,h)^{|w|}\tr(M_{gh}^w)$
    for every nonempty word $w$.
    Lemma~\ref{lem:asym_mpo_trace_eq} converts this trace identity into
    the stated operator identity for every $L\geq1$.
    The odd-ring anticommutation follows by substituting
    $\lambda(x,y)=-1$ and $\lambda(y,x)=+1$ and using $(-1)^L=-1$.
\end{proof}

\begin{theorem}[The condensation defect squares to a length-dependent
        multiple of itself]%
    \label{thm:asymex_z2z2_defect_sq}
    \lean{Z2Z2Condensation.mpo_condensationTensor_sq,
        Z2Z2Condensation.mpo_condensationTensor_apply_zero,
        Z2Z2Condensation.mpo_condensationTensor_ne_zero,
        Z2Z2Condensation.condensation_quarter_sq_eq_self_iff_even}
    \leanok
    \uses{def:asymex_z2z2_defect, thm:asymex_z2z2_fusion}
    For every $L\geq1$, $\rho^{(L)}(N)^2=(3+(-1)^L)\rho^{(L)}(N)$,
    and $\rho^{(L)}(N)\neq0$. Consequently
    $(\tfrac14\rho^{(L)}(N))^2=\tfrac14\rho^{(L)}(N)$
    if and only if $L$ is even. Thus the normalization by $\tfrac14$
    is not idempotent on odd rings.
\end{theorem}

\begin{proof}
    \leanok
    \uses{thm:asymex_z2z2_fusion}
    Expanding the square by Theorem~\ref{thm:asymex_z2z2_fusion}, each of
    the four operators $\rho^{(L)}(M_k)$ receives the four weights
    $\lambda(g,h)^L$ from the four pairs $(g,h)$ with $gh=k$, three of
    which are $+1$ and one of which is $(-1)^L$, giving the stated
    coefficient. The all-zero diagonal entry of $\rho^{(L)}(N)$ equals
    $1$, since only the identity tensor contributes to it, so
    $\rho^{(L)}(N)\neq0$. The idempotence equivalence follows by comparing
    the coefficient $\tfrac1{16}(3+(-1)^L)$ with $\tfrac14$
    on this nonzero operator.
\end{proof}

\section{The sixteen pair blocks of the stacked defect}%
\label{sec:asymex_z2z2_pairs}

The stacked product $B=N\cdot N$ has bond dimension thirty-six.
After grouping the bond indices by
$(g,h)\in(\Z_2\times\Z_2)^2$, it is block diagonal with sixteen blocks.
The block indexed by $(g,h)$ is $M_g\cdot M_h$, the stacked tensor
generating $U_gU_h$.

\begin{center}
    \tenkzkernel
    % Formula: B = N.N, stacked as two periodic rings of N contracted
    % against the sum \sum_m N^{(o,m)} (x) N^{(m,i)}, equal to one
    % periodic ring of the bond-36 tensor B (Definition~\ref{def:mpdo_operator_product}).
    % Ink/index correspondence: the two rows share the virtual bond of the
    % top ring with the physical-down legs of the bottom ring stacked
    % against the physical-up legs (contracted middle index m, drawn as
    % the shared column). Boundary signature of each side: (0,0,L,L).
    \begin{tenkz}[rows={op,op}, cols=4, boundary=periodic]
        \tn[skin=mpo, ports={90:physical:$o_1$, 270:physical}]{N} &
        \tn[skin=mpo, ports={90:physical:$o_2$, 270:physical}]{N} &
        \tn[skin=dots, ports={180:virtual, 0:virtual}]{} &
        \tn[skin=mpo, ports={90:physical:$o_L$, 270:physical}]{N} \\
        \tn[skin=mpo, ports={90:physical, 270:physical:$i_1$}]{N} &
        \tn[skin=mpo, ports={90:physical, 270:physical:$i_2$}]{N} &
        \tn[skin=dots, ports={180:virtual, 0:virtual}]{} &
        \tn[skin=mpo, ports={90:physical, 270:physical:$i_L$}]{N}
    \end{tenkz}
    \quad $=$ \quad
    \begin{tenkz}[rows={op}, cols=4, boundary=periodic]
        \tn[skin=mpo, ports={90:physical:$o_1$, 270:physical:$i_1$}]{B} &
        \tn[skin=mpo, ports={90:physical:$o_2$, 270:physical:$i_2$}]{B} &
        \tn[skin=dots, ports={180:virtual, 0:virtual}]{} &
        \tn[skin=mpo, ports={90:physical:$o_L$, 270:physical:$i_L$}]{B}
    \end{tenkz}
\end{center}

\begin{definition}[The stacked defect tensor and its sixteen target tensors]%
    \label{def:asymex_z2z2_pairs}
    \lean{Z2Z2Condensation.condensationSquare, Z2Z2Condensation.pairSlots,
        Z2Z2Condensation.pairDim, Z2Z2Condensation.pairTarget}
    \leanok
    \uses{def:asymex_z2z2_defect, def:mpdo_operator_product}
    Pairing the physical indices of $N\cdot N$ gives an MPS tensor
    on $16$ letters. Its sixteen target tensors are
    $C_{(g,h)}=\lambda(g,h)M_{gh}$, indexed by the pairs $(g,h)$.
    Their bond dimensions are those of $M_{gh}$.
\end{definition}

For twelve of the sixteen pairs, the stacked tensor is conjugate to
its target by an invertible integer matrix, with no additional zero
diagonal blocks. These are the split pairs.
Ten of them, with at most one factor from $\{x,xy\}$ and with
$\lambda(g,h)=+1$, are equal to their targets.
The two pairs $(x,y)$ and $(xy,y)$ have weight $-1$ and are conjugate
to $-M_{xy}$ and $-M_x$, respectively, by
$\begin{pmatrix}0&-1\\1&0\end{pmatrix}$.
In the following diagram, $B$ denotes one such pair block and
$s=(g,h)$.

\begin{center}
    \tenkzkernel
    % Formula: W_s B^{oi} V_s = C_s^{oi} = lambda(g,h) M_{gh}^{oi}, the
    % exact (zero-remainder) compression pair of a split pair block.
    % Ink/index correspondence: the rings W_s, V_s are the compression
    % maps; contracting them against B along the shared virtual bond and
    % against C_s on the right-hand side leaves the same free physical
    % legs o, i on both sides. Boundary signature both sides: (0,0,1,1).
    \begin{tenkz}[rows={wire}, cols=3, west=open, east=open]
        \tn[skin=ring]{W_s} &
        \tn[skin=mpo, ports={90:physical:$o$, 270:physical:$i$}]{B} &
        \tn[skin=ring]{V_s}
    \end{tenkz}
    \quad $=$ \quad
    \begin{tenkz}[rows={wire}, cols=1, west=open, east=open]
        \tn[skin=mpo, ports={90:physical:$o$, 270:physical:$i$}]{C_s}
    \end{tenkz}
\end{center}

\begin{theorem}[The twelve split pair blocks]%
    \label{thm:asymex_z2z2_split}
    \lean{Z2Z2Condensation.eEStacked_eq, Z2Z2Condensation.eYStacked_eq,
        Z2Z2Condensation.eXStacked_eq, Z2Z2Condensation.eXyStacked_eq,
        Z2Z2Condensation.yEStacked_eq, Z2Z2Condensation.yYStacked_eq,
        Z2Z2Condensation.yXStacked_eq, Z2Z2Condensation.yXyStacked_eq,
        Z2Z2Condensation.xEStacked_eq, Z2Z2Condensation.xyEStacked_eq,
        Z2Z2Condensation.xYCompression, Z2Z2Condensation.xyYCompression,
        Z2Z2Condensation.eE_remainder_eq_zero,
        Z2Z2Condensation.eY_remainder_eq_zero,
        Z2Z2Condensation.eX_remainder_eq_zero,
        Z2Z2Condensation.eXy_remainder_eq_zero,
        Z2Z2Condensation.yE_remainder_eq_zero,
        Z2Z2Condensation.yY_remainder_eq_zero,
        Z2Z2Condensation.yX_remainder_eq_zero,
        Z2Z2Condensation.yXy_remainder_eq_zero,
        Z2Z2Condensation.xE_remainder_eq_zero,
        Z2Z2Condensation.xyE_remainder_eq_zero,
        Z2Z2Condensation.xY_remainder_eq_zero,
        Z2Z2Condensation.xyY_remainder_eq_zero}
    \leanok
    \uses{def:asymex_z2z2_pairs, def:asymex_one_slot, thm:asymex_explicit_gauge}
    For each of the twelve split pairs $(g,h)$, the corresponding block
    of the stacked defect tensor admits a compression onto the single
    target $\lambda(g,h)M_{gh}$ with $z=0$.
    An invertible bond-space gauge conjugates the block exactly to this
    target, with vanishing remainder. The gauge is the identity for ten
    pairs and $\begin{pmatrix}0&-1\\1&0\end{pmatrix}$ for $(x,y)$ and
    $(xy,y)$.
\end{theorem}

\begin{proof}
    \leanok
    \uses{def:asymex_one_slot, thm:asymex_one_slot, thm:asymex_explicit_gauge}
    A finite calculation shows that ten of the blocks equal their
    targets at every physical-index pair.
    The remaining two are conjugate to their targets by the stated
    gauge, by Theorem~\ref{thm:asymex_explicit_gauge}.
    The gauge acts on the entire bond space, so the resulting
    compression has no zero diagonal blocks and its remainder is
    identically zero.
\end{proof}

\begin{theorem}[The four non-split pair blocks]%
    \label{thm:asymex_z2z2_nonsplit}
    \lean{Z2Z2Condensation.xXCompression, Z2Z2Condensation.xXyCompression,
        Z2Z2Condensation.xyXCompression, Z2Z2Condensation.xyXyCompression,
        Z2Z2Condensation.xX_evalWord_remainder_eq_zero,
        Z2Z2Condensation.xXy_evalWord_remainder_eq_zero,
        Z2Z2Condensation.xyX_evalWord_remainder_eq_zero,
        Z2Z2Condensation.xyXy_evalWord_remainder_eq_zero}
    \leanok
    \uses{def:asymex_z2z2_pairs, def:asymex_one_slot, thm:asymex_explicit_gauge,
        thm:asym_multiblock}
    The four pairs $(g,h)$ with both $g,h\in\{x,xy\}$ do not split.
    Each has bond dimension four and admits a compression onto its
    one-dimensional target $\lambda(g,h)M_{gh}$ with $z=3$ zero
    diagonal entries, an explicit unimodular integer gauge, and a
    remainder that vanishes on every word of length at least four.
    This is the bound $1+z=4$ of
    Theorem~\ref{thm:asym_multiblock}.
    For the block $(x,x)$, whose target is
    $\lambda(x,x)M_e=\delta$, the compression pair is
    \begin{align}
        W & =(0,0,0,1), \notag\\
        V & =(1,-1,-1,1)^{\mathsf T},
        \notag
    \end{align}
    with $WV=1$ and $WB^{ij}V=\delta^{ij}$ for every physical-index
    pair. The remainder $R^{ij}=B^{ij}-V\delta^{ij}W$ satisfies
    $R^{i_1j_1}R^{i_2j_2}R^{i_3j_3}R^{i_4j_4}=0$.
    The other three blocks admit the same construction, each with its
    own gauge.
\end{theorem}

\begin{proof}
    \leanok
    \uses{thm:asymex_explicit_gauge}
    For each of the four pairs, a finite calculation verifies an
    integer gauge and its inverse. This gauge puts every matrix of
    the stacked tensor in upper triangular form, with the target
    in the third diagonal position and zero in the other diagonal
    positions.
    Theorem~\ref{thm:asymex_explicit_gauge} gives the compression pair
    from this gauge. The nilpotency clause of
    Theorem~\ref{thm:asym_multiblock}, with $|S|+z=4$, gives the stated
    vanishing. The absence of a splitting follows from the
    intertwiner calculation in
    Theorem~\ref{thm:asymex_z2z2_no_intertwiner}.
\end{proof}

\begin{theorem}[Both sitewise intertwiner spaces vanish for the four
        non-split blocks]%
    \label{thm:asymex_z2z2_no_intertwiner}
    \lean{Z2Z2Condensation.xX_right_intertwiner_eq_zero,
        Z2Z2Condensation.xX_left_intertwiner_eq_zero,
        Z2Z2Condensation.xXy_right_intertwiner_eq_zero,
        Z2Z2Condensation.xXy_left_intertwiner_eq_zero,
        Z2Z2Condensation.xyX_right_intertwiner_eq_zero,
        Z2Z2Condensation.xyX_left_intertwiner_eq_zero,
        Z2Z2Condensation.xyXy_right_intertwiner_eq_zero,
        Z2Z2Condensation.xyXy_left_intertwiner_eq_zero}
    \leanok
    \uses{thm:asymex_z2z2_nonsplit}
    For each of the four non-split pairs $(g,h)$, let $B$ denote its
    stacked block and set $C^{oi}=\lambda(g,h)(M_{gh}^{oi})_{00}$.
    If a column vector $v\in\C^4$ satisfies $B^{oi}v=C^{oi}v$ for
    every physical-index pair, then $v=0$.
    Likewise, if a row vector $u$ satisfies $uB^{oi}=C^{oi}u$ for
    every physical-index pair, then $u=0$.
\end{theorem}

\begin{proof}
    \leanok
    Four of the sixteen matrices of each block are nonzero, taking
    only two distinct values. Each distinct value has rank one, with
    an explicit integer kernel and image.
    Comparing the four coordinates of the two independent equations
    forces every coordinate of $v$, respectively of $u$, to vanish.
    For the block $(x,x)$, the nonzero physical-index pairs are
    $(0,0)$ and $(1,1)$, repeated at $(2,2)$ and $(3,3)$.
    Comparing the coordinates of $B^{ij}v=v$ at $(0,0)$ gives
    $v_0=v_1=v_2=v_3$, and at $(1,1)$ gives $v_3=v_0$ and
    $v_1=-v_3$, forcing $v_0=v_1=v_2=v_3=0$.
    The row-vector comparison at $(0,0)$ gives $u_1=u_2=u_3=0$,
    and at $(1,1)$ gives $u_0=0$.
    The other three blocks eliminate coordinates in the same way,
    using their respective rank-one matrices.
\end{proof}

\begin{remark}[Failure of closure under conjugate transposition]%
    \label{rem:asymex_z2z2_anomaly}
    \uses{cor:asym_anomaly_not_star, thm:asymex_z2z2_no_intertwiner}
    By Theorem~\ref{thm:asymex_z2z2_no_intertwiner} and
    Corollary~\ref{cor:asym_anomaly_not_star}, the algebra generated
    by the matrices of each stacked pair block with
    $g,h\in\{x,xy\}$ is not closed under conjugate transposition.
    This is the same algebraic obstruction as for the squared CZX
    tensor in Section~\ref{sec:asymex_czx}.
    For each of these four blocks, there is no nonzero sitewise
    intertwiner with the specified target in either direction.
    In contrast, the twelve split pairs of
    Theorem~\ref{thm:asymex_z2z2_split} have one-dimensional
    sitewise intertwiner spaces in both directions.
    These conclusions concern the pair blocks and their target
    tensors, not the three-cocycle associator.
\end{remark}

\begin{theorem}[Compression of the stacked defect tensor]%
    \label{thm:asymex_z2z2_assembled}
    \lean{Z2Z2Condensation.condensationSquare_trace_evalWord,
        Z2Z2Condensation.exists_condensationSquare_compression,
        Z2Z2Condensation.condensationSquare_compression_z}
    \leanok
    \uses{def:asymex_z2z2_pairs, thm:asym_multiblock,
        thm:asymex_z2z2_split, thm:asymex_z2z2_nonsplit}
    For every nonempty word $w$, the stacked defect tensor $B=N\cdot N$
    satisfies
    \begin{align}
        \tr(B^w)
        &=\sum_{g,h\in\Z_2\times\Z_2}
            \lambda(g,h)^{|w|}\tr(M_{gh}^w).
        \notag
    \end{align}
    Consequently $B$ admits a multi-block compression, in the sense of
    Theorem~\ref{thm:asym_multiblock}, onto the sixteen weighted
    targets $\lambda(g,h)M_{gh}$.
    Every such compression has exactly $z=12$ zero diagonal entries,
    since $36=4(1+2+1+2)+12$.
    In the pair-block construction, these twelve entries come from
    the four non-split blocks of
    Theorem~\ref{thm:asymex_z2z2_nonsplit}, each of which has $z=3$.
\end{theorem}

\begin{proof}
    \leanok
    \uses{thm:asym_multiblock}
    Sum the word-trace identities underlying
    Theorem~\ref{thm:asymex_z2z2_fusion} over the sixteen pairs.
    Together with Lemma~\ref{lem:asym_mpo_trace_eq}, this identifies
    the word traces of the stacked defect tensor with the stated
    sum.
    Theorem~\ref{thm:asym_multiblock} then supplies the compression.
    The dimension count follows from clause (vii) and the bond
    dimensions in Theorems~\ref{thm:asymex_z2z2_split}
    and~\ref{thm:asymex_z2z2_nonsplit}.
\end{proof}

\begin{remark}[Scope of the compression results]%
    \label{rem:asymex_z2z2_scope}
    The periodic fusion table, the length-dependent identity for the
    condensation defect, and the compression results do not determine
    an element of $H^3(\Z_2\times\Z_2,U(1))$.
    The four non-split pair blocks have remainders that vanish on
    every word of length at least four; the twelve split pair blocks
    have identically vanishing remainders.
    The bound four is not asserted to be optimal, and these
    nilpotent remainders are not identified with a cohomological
    invariant.
    % The class attached to this presentation is external data in
    % Notes/OpenProblemsTN/checks/asym_z2z2_typeIII_data.md.
\end{remark}
